
\documentclass[12pt,letterpaper]{article}

\usepackage[utf8]{inputenc}
\usepackage[spanish]{babel}
\usepackage[T1]{fontenc}

\usepackage[
    top=3cm,
    left=4cm,
    bottom=3cm,
    right=2cm
]{geometry}

\usepackage{setspace}
\onehalfspacing

\usepackage{mathptmx} 

\usepackage{fancyhdr}
\pagestyle{fancy}
\fancyhf{}
\fancyfoot{}
\fancyhead[R]{\thepage}
\renewcommand{\headrulewidth}{0pt}

\usepackage{graphicx}
\usepackage{caption}
\captionsetup{font=small, labelfont=bf}

\usepackage[backend=biber, style=numeric, sorting=none]{biblatex}
\addbibresource{references.bib}   

\usepackage{titlesec}
\titleformat{\section}{\normalfont\bfseries\large}{\thesection.}{1em}{\MakeUppercase}
\titleformat{\subsection}{\normalfont\bfseries}{\thesubsection}{1em}{\MakeUppercase}

\begin{document}
\begin{titlepage}
    \thispagestyle{empty}
    \centering
    \vspace*{2cm}
    {\LARGE \textbf{CHATBOTS CONVERSACIONALES VS. FORMULARIOS TRADICIONALES EN LA UPTC}}\\[3cm]
    {\large Edwin David Fino Velandia\\ Edward Mauricio Guatibonza Suárez\\ Fredy Alejandro Bonilla Rodríguez}\\[2cm]
    {\large Universidad Pedagógica y Tecnológica de Colombia (UPTC)}\\
    {\large Facultad de Ingeniería}\\
    {\large Escuela de Ingeniería de Sistemas y Computación}\\[2cm]
    {\large Tunja}\\
    {\large \the\year}
\end{titlepage}

\thispagestyle{empty}
\vspace*{4cm}
\noindent \textbf{Nota de Aceptación:}\\[2cm]
\begin{flushright}
    \rule{6cm}{0.4pt}\\
    \rule{6cm}{0.4pt}\\
    \rule{6cm}{0.4pt}\\[1cm]
    \rule{6cm}{0.4pt}\\
    Firma del Jurado\\[1cm]
    \rule{6cm}{0.4pt}\\
    Firma del Jurado
\end{flushright}
\newpage

\thispagestyle{empty}
\tableofcontents
\newpage

\setcounter{page}{1}

\section{Introducción}
Texto aquí...

\section{Planteamiento del Problema}
\subsection{Descripción del Problema}
Texto aquí...

\subsection{Formulación del Problema}
Texto aquí...

\section{Justificación}
Texto aquí...

\section{Objetivos}
\subsection{Objetivo General}
Texto aquí...

\subsection{Objetivos Específicos}
Texto aquí...

\section{Marco de Referencia}
\subsection{Antecedentes}

La tecnología actual ha alcanzado un nivel de desarrollo que permite imitar aspectos de la inteligencia humana mediante lo que conocemos como Inteligencia Artificial (IA). Estos sistemas mantienen conversaciones coherentes y contextualizadas, resuelven problemas de distinta naturaleza, aprenden a partir de grandes volúmenes de datos y toman decisiones basadas en patrones previamente identificados. Este avance es resultado de décadas de investigación en aprendizaje automático, procesamiento del lenguaje natural y redes neuronales.
 
En el campo de la programación, la evidencia empírica confirma que este impacto no es meramente anecdótico. El experimento controlado de Peng, Kalliamvakou, Cihon y Dohmke~\cite{peng2023impact} demostró que los desarrolladores asistidos por IA completaron una tarea de referencia (la creación de un servidor HTTP en JavaScript) un 55.8\% más rápido que el grupo de control, y --dato central para este proyecto-- que los programadores con menos experiencia fueron quienes obtuvieron el mayor beneficio relativo en velocidad. Esto sugiere que el impacto de estas herramientas podría ser incluso más pronunciado en poblaciones estudiantiles que en profesionales consolidados, lo que justifica que la pregunta de rentabilidad se plantee específicamente en el contexto universitario y no solo en el laboral.
 
A nivel de uso continuado y de pago, Cui et al.~\cite{cui2024productivity} encontraron, en un experimento de campo con 4,867 desarrolladores en entornos empresariales, un aumento medio del 26.08\% en tareas completadas al usar herramientas de suscripción. Este dato empírico permite establecer un punto de comparación cuantitativo entre inversión económica y productividad neta, que es precisamente la relación costo-beneficio que este proyecto busca trasladar y verificar en el contexto académico, donde las condiciones de uso, motivación y complejidad de las tareas difieren de las de un entorno corporativo.
 
Sin embargo, la productividad medida en velocidad de ejecución no es el único eje relevante. Prather et al.~\cite{prather2023programming} argumentan que la incorporación de generadores de código basados en LLM está transformando el enfoque pedagógico de la programación: el aprendizaje se desplaza desde la memorización de sintaxis hacia la formulación de instrucciones (prompting), la depuración y la verificación lógica del código generado. Esto plantea una tensión que el proyecto debe considerar: ganar tiempo no equivale necesariamente a aprender mejor, y por tanto la rentabilidad de una suscripción de pago no puede medirse solo en tiempo ahorrado, sino también en si esa ganancia de tiempo se traduce en comprensión técnica real o en dependencia.
 
Esta tensión se ve reforzada por los hallazgos de Martinović y Rozić~\cite{martinovic2024ai}, quienes muestran que, si bien las herramientas de IA reducen errores sintácticos iniciales, un uso desmedido y sin criterio de revisión puede derivar en código menos estructurado o en vacíos de comprensión lógica. Este hallazgo es clave para cuestionar si la versión gratuita --con sus limitaciones-- es en realidad suficiente, o incluso preferible, durante la fase de fundamentación académica, precisamente porque impone restricciones que obligan a un mayor esfuerzo cognitivo del estudiante.
 
La revisión sistemática de más de 60 estudios empíricos sobre IA generativa en la educación en ciencias de la computación (2025/2026)~\cite{dialnet2025inteligencia,comillas2025impacto,ups2025inteligencia,scielo2025revision,redalyc2024revision,scielochile2025uso,horizontes2024implicaciones,zenodo2025impacto} aporta una base metodológica más amplia: evalúa la carga cognitiva del estudiante y compara las tasas de alucinación de código entre modelos gratuitos y modelos avanzados de pago, lo que da sustento sistemático a la hipótesis de que invertir en versiones avanzadas podría reducir el tiempo de depuración --uno de los indicadores centrales que este proyecto se propone medir--.
 
Finalmente, el marco de la ``Productivity-Validation Tension'' descrito por IEEE Software~\cite{ieee2024assisted} resulta metodológicamente decisivo para este proyecto: aunque las herramientas de suscripción reducen el tiempo dedicado a escribir código repetitivo, introducen un costo adicional en la validación y prueba de las sugerencias generadas. Esto respalda directamente la decisión metodológica del proyecto de no medir la rentabilidad únicamente en ``líneas escritas por minuto'', sino en el tiempo neto ahorrado, considerando también el proceso de revisión y comprobación del trabajo académico. Del mismo modo, Ziegler et al.~\cite{ziegler2022measuring} muestran que tasas de aceptación de sugerencias cercanas al 27\%--30\% se correlacionan con una mayor percepción de fluidez y productividad, lo que permite contrastar el comportamiento de un usuario con acceso ilimitado frente a las restricciones propias de un plan gratuito.
 
En conjunto, esta evidencia respalda la pertinencia de evaluar si el costo de una suscripción de pago se justifica por la mejora real en productividad, calidad de código y aprendizaje autónomo del estudiante, o si la versión gratuita resulta suficiente en las etapas iniciales de formación, sin requerir esa inversión adicional.
    
\subsection{Marco Teórico}
Texto aquí...

\subsection{Marco Conceptual}
Texto aquí...

\subsection{Marco Legal}
Texto aquí... 

\section{Diseño Metodológico}
\subsection{Tipo de Investigación}
Texto aquí...

\subsection{Población y Muestra}
Texto aquí...

\subsection{Técnicas e Instrumentos de Recolección de Información}
Texto aquí...

\subsection{Fases de la Investigación}
Texto aquí...

\section{Cronograma de Actividades}
Texto aquí... 

\section{Presupuesto}
Texto aquí...

\printbibliography[title={Bibliografía}]
\newpage
\section*{Anexos}
Aquí van tablas, encuestas o documentos externos.

\end{document}

